  In diesem Kapitel geht es erstmal darum die Mathematische Basis Aufzubauen. Wie aus der Einleitung hervorgeht, wollen wir jeweils zwei Paare von Stimuli auf ihre Ähnlichkeit vergleichen. Seien $a, b, c$ und $d$ solche Stimuli. Dann bedeutet die Notation $(a,b)\succsim(c,d)$, dass Stimuli $a$ und $b$ sich unähnlicher sind, als $c$ und $d$. Analog bedeutet $\sim$, dass die Ähnlichkeit zweier Paare gleich ist.\\
  Die Idee von einer proximity structure, ist die Mathematische Definition dieser Idee: eine Relation $\succsim$ auf einer Menge $A$ die die Grundeigenschaften einer Metrik erfüllt.

  \begin{defi}[factorial proximity structure]
Suppose $A$ is a set and $\succsim$ a quaternery relation on $A$, i.e., binary relation on $A \times A$. $\langle A, \succsim \rangle$ is a proximity structure iff the following axioms hold for all $a, b \in A$:

\begin{enumerate}
  \item $\succsim$ is a weak ordering on $A \times A$, i.e., it is connected and transitive.
    \item $(a, b) \succsim (a, a)$ whenever $a \neq b$.
    \item $(a, a) \sim (b, b)$ (minimality).
    \item $(a, b) \sim (b, a)$ (symmetry).
\end{enumerate}

The structure is factorial iff
\[
A = \prod_{i=1}^{n} A_i.
\]
\end{defi}

\begin{rem}
Zusammenhängend heißt, man kann alle distinkten Elemente vergleichen, also für eine Relation $R$ auf einer Menge $X$ gilt für $x, y \in X, x \neq y$ dann $xRy$ oder $yRx$.
\end{rem}

Da wir die Relation $\succsim$ durch eine Metrik repräsentieren wollen, lohnt es sich einen Begriff dafür einzuführen. Die Idee ist das, was man sich intuitiv vorstellt: man will dass die Metrik die Struktur von $\langle A, \succsim \rangle$ erhält, und die einzige Struktur die wir gegeben haben, ist die Ordnung jeweils zweier Paare.

\begin{defi}[Representability of $\succsim$ by a metric]
A proximity structure $\langle A, \succsim \rangle$ is representable by a metric iff there exists a real-valued function $\delta$ on $A$ such that:

\begin{enumerate}
    \item $\langle A, \delta \rangle$ is a metric space, i.e., $\delta$ is a metric on $A$.
    \item $\delta$ preserves the relation $\succsim$, i.e., for all $a, b, c, d \in A$
    \[
    \delta(a, b) \geq \delta(c, d) \iff (a, b) \succsim (c, d).
    \]
\end{enumerate}
\end{defi}
